% ══════════════════════════════════════════════════════════════════════════════
\section{Conclusion}
\label{sec:conclusion}
% ══════════════════════════════════════════════════════════════════════════════

This paper estimates the effect of satellite-predicted maize yields on violent conflict across seven Sub-Saharan African countries using an 18,096 observation Admin-2~$\times$~year panel from 2010 to 2024. By using ML-predicted rather than observed yields as the key regressor---constructed entirely from exogenous environmental and satellite inputs via an XGBoost model trained on GROW-Africa administrative records---the analysis avoids the reverse-causality bias that makes observed-yield regressions uninformative about the causal direction of the agriculture--conflict relationship.

Three main findings emerge. First, higher predicted maize yields reduce conflict on average: the preferred TWFE OLS log-log specification estimates an elasticity of $-0.666$ ($p < 0.001$), which survives spatial lag control, strengthens at longer post-harvest horizons, and is concentrated in violence against civilians and battles rather than riots. This pattern is consistent with the opportunity-cost channel: agricultural productivity raises the return to peaceful production, reducing the incentive for armed actors and potential recruits to engage in violence.

Second, this aggregate pacifying effect conceals a fundamental heterogeneity in direction across agro-ecological zones. In Sahel countries---Burkina Faso, Mali, and Niger---where armed groups systematically appropriate agricultural surplus through coercive market taxation and territorial control of productive zones, higher predicted yields are associated with \textit{more} conflict. The full-sample interaction coefficient of $+3.513$ ($p < 0.001$) produces a net Sahel elasticity of $+2.383$: the same improvement in agricultural conditions that suppresses conflict in Non-Sahel environments amplifies it in Sahel ones. This finding constitutes direct evidence for the resource-predation mechanism in the Sahel, and demonstrates that the sign of the agriculture--conflict relationship is context-dependent in a way that cannot be recovered from pooled estimates.

Third, the null result for riots and the significant negative effect on violence against civilians confirm that the mechanism operates specifically at the armed-actor margin rather than through broad civilian grievance or economic discontent. Agricultural income shocks affect the profitability of predatory violence, not the propensity of civilian populations to protest.

These findings carry practical implications. In Non-Sahel environments, agricultural investments that increase productivity can be expected to reduce local conflict incidence as a co-benefit---a result that supports the case for aid-financed agricultural programmes in fragile but non-Sahelian contexts. In the Sahel, this co-benefit logic does not hold: productivity-enhancing interventions may inadvertently increase the value of agricultural surplus available for extraction, attracting armed group activity unless accompanied by complementary investments in property rights enforcement, market security, and institutional capacity to protect the gains from improved production.

Several limitations should be kept in mind. The yield model's spatial $R^2$ of 0.35 implies substantial unexplained variance, and measurement error in the predictor attenuates estimates toward zero, suggesting the true elasticities may be larger in magnitude than those reported. The Sahel classification is country-level and does not capture within-country agro-ecological heterogeneity---notably in northern Nigeria and southern Mali---which a finer-grained AEZ variable would address. The causal interpretation rests on the assumption that satellite-derived vegetation signals are not directly affected by conflict displacement at the district scale, which is plausible but not fully testable within this design.

Future work should exploit finer-grained FAO GAEZ classifications to test the predation mechanism at the sub-country level, extend the analysis to other food crops using multi-crop prediction models, and examine the distributional consequences of yield shocks for household-level recruitment into armed groups in Sahel contexts where administrative microdata is available.
